\documentclass[aps,prd,twocolumn,superscriptaddress,nofootinbib,10pt]{revtex4-2}

\usepackage{amsmath,amssymb}
\usepackage[colorlinks=true,citecolor=blue,linkcolor=blue,urlcolor=blue]{hyperref}

\begin{document}

\title{A rotating condensate tracking $V\propto r^{n}$ cannot reach kination:\\
exact equation of state $w=(n-2)/(n+2)$}

\author{Justin Pulford}
\email{pulfordj420@gmail.com}
\affiliation{Unaffiliated}

\date{\today}

\begin{abstract}
A complex scalar with a conserved charge $Q$, written in polar form as a rotating condensate of
amplitude $r$, has rotational energy density that would scale as $a^{-6}$ (kination, $w=1$) if $r$
were held fixed. On a polynomial potential $V\propto r^{n}$ the amplitude does not stay fixed: it
tracks the minimum of the effective radial potential formed by $V$ plus the centrifugal barrier.
Tracking yields the exact equation of state
\begin{equation*}
  w = \frac{n-2}{n+2}\,,
\end{equation*}
strictly below unity for every finite polynomial degree $n$, and equal to the kination limit only
as $n\to\infty$. Freeze-out of $r$ (the only escape to $w=1$) requires a trans-Planckian amplitude,
$r\gtrsim\mathcal{O}(1)\,M_{\mathrm{Pl}}$. Numerical integration of the radial field equation through a
contracting background reproduces the analytic $w$ to a few parts in $10^{5}$ at $n=2,4,6$. The
result is classical field theory with a conserved charge; it is independent of any particular
dark-sector model. It is a negative sector result: a rotating condensate of this class cannot
supply the stiff phase often invoked to survive a BKL approach to a crunch.
\end{abstract}

\maketitle

% ---------------------------------------------------------------------------
\section{Motivation}\label{sec:motivation}

In a contracting cosmology the shear energy density of anisotropic modes scales as $a^{-6}$.
Radiation ($w=1/3$) scales as $a^{-4}$, so the anisotropy-to-radiation ratio grows as $a^{-2}$
and becomes catastrophic over even a few decades of contraction. The standard route past a
Belinski--Khalatnikov--Lifshitz (BKL) approach~\cite{BKL1970,BKL1982} is a stiff phase with
$w\geq 1$ (kination), which scales as $a^{-6}$ and keeps the relative anisotropy from growing.

A rotating complex scalar looks like a natural supplier of that phase. Written
$\Psi=(r/\sqrt{2})\,e^{i\theta}$ with conserved charge
\begin{equation}
  Q = a^{3} r^{2}\dot\theta\,,
\end{equation}
elimination of $\dot\theta$ at fixed $Q$ produces a centrifugal term in the effective radial
potential,
\begin{equation}
  V_{\mathrm{eff}}(r)
  = V(r) + \frac{Q^{2}}{2 a^{6} r^{2}}\,.
  \label{eq:Veff}
\end{equation}
If the amplitude $r$ were held fixed, the rotational energy density
$\rho_{\mathrm{rot}}=Q^{2}/(2a^{6}r^{2})$ would fall exactly as $a^{-6}$ --- pure kination.
The question is whether $r$ holds fixed on approach. That is decidable from the field equation
alone.

This note answers that question for polynomial potentials $V\propto r^{n}$. The answer is
negative at every sub-Planckian amplitude. No claim is made about bounce mechanisms, dark energy,
or any model beyond the classical charged scalar.

% ---------------------------------------------------------------------------
\section{Tracking equation of state}\label{sec:tracking}

At the minimum of~\eqref{eq:Veff},
\begin{equation}
  V'(r) = \frac{Q^{2}}{a^{6} r^{3}}\,.
  \label{eq:min}
\end{equation}
For $V=c\,r^{n}$ this integrates immediately to
\begin{equation}
  r \propto a^{-6/(n+2)}\,.
  \label{eq:rtrack}
\end{equation}
The rotational kinetic energy is then related to the potential by
\begin{equation}
  \rho_{\mathrm{rot}}
  = \frac{Q^{2}}{2 a^{6} r^{2}}
  = \frac{n}{2}\,V\,,
\end{equation}
so the total density and pressure at tracking are
\begin{align}
  \rho &= V + \rho_{\mathrm{rot}} = V\bigl(1+n/2\bigr)
       = V\,(n+2)/2\,, \\
  p    &= \rho_{\mathrm{rot}} - V = V\bigl(n/2-1\bigr)
       = V\,(n-2)/2\,.
\end{align}
Hence
\begin{equation}
  \boxed{\,w = \frac{p}{\rho} = \frac{n-2}{n+2}\,}\,.
  \label{eq:w}
\end{equation}
Equivalently, from~\eqref{eq:rtrack},
$\rho\propto a^{-6n/(n+2)}=a^{-3(1+w)}$, which recovers the same $w$.

Special cases reproduce textbook regimes rather than inventing new ones:
\begin{itemize}
\item $n=2$ (quadratic / mass era): $w=0$ (cold dust / oscillating massive scalar).
\item $n=4$ (quartic): $w=1/3$ (radiation-like).
\item $n\to\infty$: $w\to 1$ (kination approached only at infinite degree).
\end{itemize}
For every finite polynomial $n$, $w<1$. No polynomial potential reaches the stiff condition by
tracking.

The algebraic relation~\eqref{eq:w} is not new in cosmology: it is the standard effective
equation of state of a coherently oscillating scalar in a monomial potential~\cite{Turner1983}
(the cold-dark-matter reading of a massive scalar is the $n=2$ case). Kination itself is the
$w=1$ end-member of kinetic domination~\cite{Spokoiny1993}. What is stated here is the
consequence for a \emph{rotating} condensate whose centrifugal barrier might have been hoped to
supply kination without freezing $r$.

% ---------------------------------------------------------------------------
\section{Why the field tracks rather than freezes}\label{sec:freeze}

Freeze-out of $r$ would leave $\rho_{\mathrm{rot}}\propto a^{-6}$ and restore $w=1$. Freeze-out
requires Hubble friction to beat the restoring curvature of $V_{\mathrm{eff}}$ at the tracking
minimum, i.e.\ $V_{\mathrm{eff}}''\lesssim H^{2}$.

At the minimum~\eqref{eq:min}, for $V=c\,r^{n}$,
\begin{equation}
  V_{\mathrm{eff}}''
  = n(n+2)\,\frac{V}{r^{2}}\,,
  \qquad
  H^{2}
  = \frac{\rho}{3M_{\mathrm{Pl}}^{2}}
  = \frac{V(n+2)}{6M_{\mathrm{Pl}}^{2}}\,,
\end{equation}
so
\begin{equation}
  \frac{V_{\mathrm{eff}}''}{H^{2}}
  = 6n\,\Bigl(\frac{M_{\mathrm{Pl}}}{r}\Bigr)^{2}\,.
  \label{eq:curv}
\end{equation}
The ratio is $\mathcal{O}(10)\,(M_{\mathrm{Pl}}/r)^{2}$. Freezing needs $V_{\mathrm{eff}}''/H^{2}\sim 1$,
hence
\begin{equation}
  r \gtrsim \sqrt{6n}\,M_{\mathrm{Pl}}\,.
\end{equation}
Numerically this is $r\gtrsim 3.5\,M_{\mathrm{Pl}}$ at $n=2$ and
$r\gtrsim 4.9\,M_{\mathrm{Pl}}$ at $n=4$.
For any \emph{sub-Planckian} amplitude the curvature dominates by orders of magnitude, the field
tracks, and $w$ is pinned at~\eqref{eq:w}. Kination from freeze-out is available only at
trans-Planckian amplitude.

% ---------------------------------------------------------------------------
\section{Numerical check}\label{sec:numerics}

The analytic result can be confirmed by integrating the radial field equation through a
contracting FLRW background and measuring $w$ from the slope of $\ln\rho$ against $\ln a$. In
e-fold time $N=\ln a$ with $u=\mathrm{d}r/\mathrm{d}N$,
\begin{equation}
  r'' + \bigl(3 + H'/H\bigr)r'
  + \frac{1}{H^{2}}\Bigl[V'(r)-\frac{Q^{2}}{a^{6}r^{3}}\Bigr] = 0\,,
\end{equation}
with $H^{2}$ recomputed each step from the Friedmann constraint. Starting on the tracking
solution and contracting over several decades of $a$, one finds (sub-Planckian window):

\begin{center}
\begin{tabular}{c c c}
\hline
$n$ & $w$ analytic & $w$ integrated \\
\hline
2 & $0$ & $0.000064$ \\
4 & $1/3$ & $0.333395$ \\
6 & $1/2$ & $0.500067$ \\
\hline
\end{tabular}
\end{center}

Agreement is at the level of a few parts in $10^{5}$. The same integrations cross the
trans-Planckian freeze threshold when $r$ grows as $a$ falls; above that threshold the measured
$w$ turns to $1.000$ (to four figures), confirming that freeze-out does produce kination and that
the price is the amplitude~\eqref{eq:curv} already prices.

% ---------------------------------------------------------------------------
\section{What this does and does not claim}\label{sec:scope}

\textbf{Does claim.} For a rotating complex scalar with conserved $Q$ tracking
$V_{\mathrm{eff}}$ on $V\propto r^{n}$, the equation of state is exactly~\eqref{eq:w}. No finite
polynomial $n$ reaches kination. Freeze-out to $w=1$ requires $r\sim\mathcal{O}(M_{\mathrm{Pl}})$.

\textbf{Does not claim.} That no stiff component can exist in a contracting universe; only that
\emph{this} sector, at sub-Planckian amplitude, cannot supply one. Other fields, non-polynomial
potentials, non-canonical kinetics, or a bounce mechanism that introduces a stiff component of its
own are outside the argument. The note does not solve the BKL problem and does not construct a
bounce.

The result is therefore a negative sector statement of the kind useful in model-building
bookkeeping: if a construction relies on a rotating condensate of this class to protect a
contracting phase against BKL growth, that reliance fails at every sub-Planckian amplitude for
every polynomial potential.

% ---------------------------------------------------------------------------
\begin{acknowledgments}
Independent numerical integration of the radial equation was used only to cross-check the analytic
tracking law; the law itself follows from~\eqref{eq:min}--\eqref{eq:w} without numerics.
\end{acknowledgments}

\begin{thebibliography}{9}

\bibitem{BKL1970}
V.~A. Belinskii, I.~M. Khalatnikov, and E.~M. Lifshitz,
Adv.\ Phys.\ \textbf{19}, 525 (1970).

\bibitem{BKL1982}
V.~A. Belinskii, I.~M. Khalatnikov, and E.~M. Lifshitz,
Adv.\ Phys.\ \textbf{31}, 639 (1982).

\bibitem{Turner1983}
M.~S. Turner,
Phys.\ Rev.\ D \textbf{28}, 1243 (1983).

\bibitem{Spokoiny1993}
B.~Spokoiny,
Phys.\ Lett.\ B \textbf{315}, 40 (1993).

\end{thebibliography}

\end{document}
